\scanpage{349}
\mspage{88}
% archive sheet number: 348
\deleted{\unsure{Considér}} Soit donc $\alpha\in[1,s=r']$, et considérons la flèche \inserted{de tr.} $F$-critique $\underline{u}_{i_\alpha}$, soit $j_\alpha$ l'indice marquant l'intervalle critique \inserted{de $I_\tau$} contenant la \unsure{partie|portée} de $\underline{u}_{i_\alpha}$ de $I_\tau$
\[
\underline{u}_{i_\alpha}\in\mathrm{Fl}(I_{j_\alpha}),
\]
alors $F$ induit des morphismes \inserted{(croissants)} d'intervalles or.
\[
F_\alpha\doteq F|I_{j_\alpha}:I_{j_\alpha}\longrightarrow I'_\alpha
\]
et est trivial sur les autres intervalles critiques (ceux pas de la forme $I_{j_\alpha}$). \deleted{\unsure{Car $I_{j_\alpha}$. D'autre part, $F_\alpha$ n'est pas}} \inserted{La donnée des $F_\alpha$ détermine donc $F$} \deleted{\unsure{par recollage}} sur les $I_{j_\alpha}$, donc sur \uline{tous} les $I_{\unsure{j|i}}$, donc partout (par l'opération $*$ \deleted{\unsure{$\}$}} de composition par \unsure{recollement|accolement}). De plus \deleted{les $F_\alpha$} la condition a)
assurera que \deleted{les} \inserted{chaque} $F_\alpha$ \deleted{\unsure{sont}} de nature \inserted{assez} particulière~: une seule flèche de transition critique $\underline{u}_{i_\alpha}$. De plus, pour qu'on ait $F=F_0$, il faudra

b) \deleted{\unsure{La flèche de transition}} \inserted{$\forall\alpha\in[1,r']$} $\underline{u}_{i_\alpha}$ est la \uline{première} flèche de transition de l'intervalle critique $I_{j_\alpha}$ qui la contient.

Et enfin

c) La suite des intervalles critiques \deleted{\unsure{$I$}} de $I_\tau$
\[
I_{j_1},I_{j_2}\ldots I_{j_{r'}}\qquad(\text{de signes }\inserted{\text{alternés}}\ \varepsilon_1,-\varepsilon_1,\varepsilon_1,\ldots)
\]
associés aux flèches $F$-critiques, n'est autre que celle \deleted{des}
\[
I_1,I_2,\ldots I_{r'}
\]
des $r'$ premiers intervalles critiques de $I_\tau$, i.e. les
\[
(j_1,j_2,\ldots,j_{r'})=(1,2,\ldots,r').
\]
